% Actual class: 02
% Slide decks: CZ3005 Module 2_Intelligent Agents and Search; CZ3005 Module 3_Constraint satisfaction and adversarial search
% Exact scope: Module 2 pp. 29--68; Module 3 pp. 1--25, aligned with the Week 3 schedule
% NTULearn supplement: None
% Human verified: Concepts discussed and checks completed

\section{第 02 节课：Search 与 Constraint Satisfaction}
\label{lecture:SC3000-L02}

\lecturemeta
  {SC3000-L02}
  {2026-08-24}
  {CZ3005 Module 2；CZ3005 Module 3}
  {Module 2 pp.~29--68；Module 3 pp.~1--25（对应 Week 3 schedule）}
  {无}
  {知识内容已讨论并通过检查题}

\subsection{本讲主线}

Search（搜索）在 state space（状态空间）中反复选择 frontier node（边界节点）并生成 successors（后继节点）；不同算法的核心差异是 node-expansion order（节点展开顺序）。Uninformed search（无信息搜索）只使用问题定义，informed search（有信息搜索）再利用 heuristic（启发信息）。Constraint Satisfaction Problem（约束满足问题，CSP）把目标改写为变量赋值，并用 backtracking、constraint propagation 与 heuristics 尽早剪除不可能分支。

\lecturetopic{SC3000-L02-T01}{General Search 与评价标准}

Frontier 保存已生成但尚未展开的 nodes；explored set（已探索集合）避免 graph search 中重复展开相同 state。Search strategy 决定从 frontier 取出哪个 node，通常比较：

\begin{itemize}
  \item Completeness（完备性）：存在解时是否一定找到；
  \item Optimality（最优性）：是否保证 minimum-cost solution（最小代价解）；
  \item Time complexity（时间复杂度）：生成或展开的节点数；
  \item Space complexity（空间复杂度）：同时保存在内存中的最大节点数。
\end{itemize}

设 $b$ 为 maximum branching factor（最大分支因子），$d$ 为 shallowest solution depth（最浅解深度），$m$ 为 state-space maximum depth，$\ell$ 为 depth limit。Node 不仅包含 state，还需要 parent、incoming action、depth 与 path cost，才能恢复 solution path 并进行 cost-based ordering（按代价排序）。

\begin{figure}[htbp]
  \centering
  \resizebox{0.94\textwidth}{!}{\begin{tikzpicture}[
  >=Latex,
  every node/.style={font=\small},
  source/.style={draw=noteBlue,rounded corners,thick,minimum width=53mm,minimum height=34mm,align=left,inner sep=3mm},
  rule/.style={draw=ntuRed,rounded corners,thick,minimum width=48mm,minimum height=10mm,align=left}
]
  \node[source] (frontier) {\textbf{Frontier}\\[1mm]
    $A:(d=2,g=5,h=1,f=6)$\\
    $B:(d=3,g=3,h=5,f=8)$\\
    $C:(d=1,g=7,h=0,f=7)$};

  \node[rule,right=18mm of frontier,yshift=18mm] (bfs) {BFS: minimum $d\Rightarrow C$};
  \node[rule,below=4mm of bfs] (ucs) {UCS: minimum $g\Rightarrow B$};
  \node[rule,below=4mm of ucs] (greedy) {Greedy: minimum $h\Rightarrow C$};
  \node[rule,below=4mm of greedy] (astar) {A*: minimum $g+h\Rightarrow A$};

  \draw[->,thick] (frontier.east) -- (bfs.west);
  \draw[->,thick] (frontier.east) -- (ucs.west);
  \draw[->,thick] (frontier.east) -- (greedy.west);
  \draw[->,thick] (frontier.east) -- (astar.west);
\end{tikzpicture}
}
  \caption{同一 frontier 在不同 priority rule（优先规则）下产生不同的下一展开节点。}
  \label{fig:sc3000-l02-search-priorities}
\end{figure}

\lecturetopic{SC3000-L02-T02}{Uninformed Search Strategies}

\begin{center}
\small
\begin{tabularx}{\textwidth}{@{}>{\raggedright\arraybackslash}p{2.1cm}>{\raggedright\arraybackslash}p{3.1cm}>{\centering\arraybackslash}p{1.7cm}>{\centering\arraybackslash}p{1.7cm}X@{}}
  \toprule
  Algorithm & Frontier rule & Time & Space & Guarantee（常用假设下） \\
  \midrule
  BFS & minimum depth；FIFO & $O(b^d)$ & $O(b^d)$ & Complete；equal step cost 时 optimal \\
  UCS & minimum $g(n)$；priority queue & $N(g\le C^*)$ & 同阶 & $c\ge\varepsilon>0$ 时 complete、optimal \\
  DFS & maximum depth；LIFO & $O(b^m)$ & $O(bm)$ & 无限深或有环时不 complete；不 optimal \\
  DLS & DFS with cutoff $\ell$ & $O(b^\ell)$ & $O(b\ell)$ & 仅当 $d\le\ell$ 才 complete；不 optimal \\
  IDS & 依次运行 DLS$(0,1,\ldots)$ & $O(b^d)$ & $O(bd)$ & Complete；equal step cost 时 optimal \\
  \bottomrule
\end{tabularx}
\end{center}

Breadth-First Search（广度优先搜索，BFS）按 depth 搜索，不比较实际 cost。Uniform-Cost Search（统一代价搜索，UCS）按
\[
  g(n)=\text{cost from initial state to }n
\]
排序；BFS 是 unit step cost 下的 UCS。UCS 不能在 goal 刚被生成时停止，必须等 goal 以 frontier 中最小的 $g$ 被取出。

Depth-First Search（深度优先搜索，DFS）内存较省，但可能陷入无限路径或 cycle。Depth-Limited Search（深度限制搜索，DLS）用 $\ell$ 截断，Iterative Deepening Search（迭代加深搜索，IDS）则逐步增大 $\ell$，以重复少量浅层工作换取 BFS 的保证与 DFS 的空间优势。

Bidirectional Search（双向搜索）若能从明确 goal 反向生成 predecessors，并能检测两侧相遇，可把规模从约 $b^d$ 降至 $b^{d/2}$；代价是两侧 frontier 仍需大量内存。

\textbf{对应例题：}\relatedexercise{SC3000-L02-E01}{Romania Route Finding：Greedy 与 A*}。

\lecturetopic{SC3000-L02-T03}{Heuristic、Greedy 与 A*}

Heuristic function（启发函数）
\[
  h(n)=\text{estimated cheapest cost from }n\text{ to a goal}
\]
使用 problem-specific knowledge（问题特定知识）改变展开顺序。Greedy Best-First Search（贪心最佳优先搜索）只比较 $h(n)$，通常搜索较快，但一般不 complete、也不 optimal。

A* 同时考虑已付代价与剩余估计：
\[
  \boxed{f(n)=g(n)+h(n)}.
\]
当 $h=0$ 时退化为 UCS；忽略 $g$ 时退化为 Greedy。A* 的关键终止条件是：\emph{goal 以 frontier 中最小的 $f$ 被取出}，而不是 goal 第一次被生成。

Admissible heuristic（可采纳启发函数）满足
\[
  0\le h(n)\le h^*(n),
\]
即从不高估真实剩余代价。Graph search 若不重新打开 closed node，通常还要求 consistency（一致性）：
\[
  h(n)\le c(n,n')+h(n').
\]
在 finite branching 与 positive lower-bounded step cost 等标准条件下，这些性质给出 A* 的 completeness 与 optimality。好的 heuristic 能显著减少展开量，但 A* 往往仍受 memory consumption 限制。

Module 2 最后用 AlphaGo 说明巨大 game tree 的搜索。标准术语是 Monte Carlo Tree Search（蒙特卡洛树搜索，MCTS），包含 selection、expansion、simulation 与 backpropagation；policy/value networks 用于引导和估值。它与 $f=g+h$ 只有概念类比，并不等同于普通 A*。

\textbf{对应例题：}\relatedexercise{SC3000-L02-E01}{Romania Route Finding：Greedy 与 A*}。

\lecturetopic{SC3000-L02-T04}{CSP Formulation 与 Backtracking}

Constraint Satisfaction Problem 可写成
\[
  \boxed{\mathrm{CSP}=(X,D,C)},
\]
其中 $X=\{X_1,\ldots,X_n\}$ 是 variables（变量），$D_i$ 是 $X_i$ 的 domain（取值域），$C$ 是 constraints（约束）。Consistent assignment 不违反已能检查的 constraints；solution 必须 complete 且 consistent。

以 4-Queens 为例，令 $Q_i$ 表示第 $i$ 行皇后所在列，$Q_i\in\{1,2,3,4\}$。对任意 $i<j$：
\[
  \boxed{Q_i\ne Q_j},\qquad
  \boxed{|Q_i-Q_j|\ne |i-j|},
\]
分别排除同列与同对角线；$i\ne j$ 只说明行号不同，不能替代 $Q_i\ne Q_j$。

Backtracking search（回溯搜索）每次选择一个未赋值 variable、尝试一个 value，并在 partial assignment 已不 consistent 时立即撤销。剪枝是安全的，因为已经违反 constraint 的 partial assignment 不可能通过增加赋值重新变合法。

Constraint propagation（约束传播）进一步把一次赋值的影响传给其他 domains；若某个 domain 变为 $\varnothing$，当前分支可立即失败。Forward checking（前向检查）只删除与新赋值直接冲突的邻居取值，更强的 propagation 可重复传播由 domain 缩减产生的新约束。

\begin{figure}[htbp]
  \centering
  \resizebox{0.96\textwidth}{!}{\begin{tikzpicture}[
  >=Latex,
  every node/.style={font=\small},
  flowstep/.style={draw=noteBlue,rounded corners,thick,minimum width=29mm,minimum height=11mm,align=center},
  decision/.style={draw=ntuRed,rounded corners,thick,minimum width=31mm,minimum height=11mm,align=center}
]
  \node[flowstep] (variable) {Select variable};
  \node[flowstep,right=10mm of variable] (value) {Order values};
  \node[flowstep,right=10mm of value] (assign) {Assign one value};
  \node[flowstep,right=10mm of assign] (propagate) {Propagate\\constraints};
  \node[decision,right=10mm of propagate] (check) {Consistent?\\Complete?};

  \node[flowstep,above=11mm of check] (solution) {Return solution};
  \node[flowstep,below=11mm of assign] (undo) {Undo assignment\\(backtrack)};

  \draw[->,thick] (variable) -- (value);
  \draw[->,thick] (value) -- (assign);
  \draw[->,thick] (assign) -- (propagate);
  \draw[->,thick] (propagate) -- (check);
  \draw[->,thick] (check) -- node[right] {complete} (solution);
  \draw[->,thick] (check.south) |- node[pos=0.2,right] {conflict / $D_i=\varnothing$} (undo.east);
  \draw[->,thick] (undo.west) -| node[pos=0.25,below] {try next value} (value.south);
  \draw[->,thick] (check.north west) |- node[pos=0.25,above] {consistent, incomplete} (variable.north);
\end{tikzpicture}
}
  \caption{Backtracking 与 constraint propagation 的控制流程。}
  \label{fig:sc3000-l02-csp-flow}
\end{figure}

\textbf{对应例题：}\relatedexercise{SC3000-L02-E02}{4-Queens 的 CSP 约束与剪枝}。

\lecturetopic{SC3000-L02-T05}{CSP Heuristics 与 Min-Conflicts}

\begin{center}
\small
\begin{tabularx}{\textwidth}{@{}>{\raggedright\arraybackslash}p{4.0cm}XX@{}}
  \toprule
  Heuristic & 选择规则 & 目的 \\
  \midrule
  Most Constraining Variable（约束最多变量） & 选择与最多 unassigned variables 相连的变量 & 尽早影响更多 domains、暴露冲突 \\
  Least Constraining Value（最少约束值，LCV） & 选择排除邻居合法取值最少的 value & 为后续赋值保留最大 flexibility \\
  Min-Conflicts（最小冲突） & 从 complete assignment 出发，修改冲突变量使 violated constraints 最少 & 用 local repair 快速处理大型 CSP \\
  \bottomrule
\end{tabularx}
\end{center}

讲义的 Most Constraining Variable 是 degree heuristic（度启发式）；不要与 Most Constrained Variable/MRV（最少剩余值）混淆，后者选择当前合法取值最少的变量。Min-Conflicts 属于 local search（局部搜索），可能遇到 plateau 或循环，通常不提供 completeness 保证。

\textbf{一分钟回顾。}BFS 比较 depth，UCS 比较 $g$，Greedy 比较 $h$，A* 比较 $g+h$；A* 在 goal 成为 frontier 最小项时停止。CSP 用 variables、domains 与 constraints 表示赋值问题；backtracking 撤销失败选择，constraint propagation 主动缩小 domains。Variable heuristic 决定先选谁，value heuristic 决定先试什么值，Min-Conflicts 从完整但冲突的赋值进行局部修复。
